% !PRINTED 301
\weberpage{301}

Suppose then that for a value $\xi$ of $x$ between $\alpha$ and $\beta$, in the sequence

\weq{3}{f^{(\nu)}(x),\ f^{(\nu+1)}(x),\ .\ .\ .\ f^{(\nu+\mu-1)}(x),\ f^{(\nu+\mu)}(x)}

all the terms, with the exception of the last, vanish, and suppose the last, $f^{(\nu+\mu)}(x)$, has, say, a positive value.

We mark off about $\xi$ two intervals $\delta_1$, $\delta_2$, such that all values of $x$ in the interval $\delta_1$ are smaller than $\xi$, and in the interval $\delta_2$ greater, and we take these intervals so small that the functions (3) do not vanish in them except at $\xi$, so that in particular $f^{(\nu+\mu)}(x)$ remains positive.

Now since, when $f^{(\nu+\mu)}(x)$ is positive, $f^{(\nu+\mu-1)}(x)$ increases together with $x$, we have

\[
\begin{array}{l}
f^{(\nu+\mu-1)}(x) \text{ in } \delta_1 \text{ negative,} \\
\text{in } \delta_2 \text{ positive.}
\end{array}
\]

From this it follows that $f^{(\nu+\mu-2)}(x)$ decreases in $\delta_1$ and increases in $\delta_2$, hence is positive in both intervals; and continuing in this way we infer the succession of signs in the sequence (3):

\[
\begin{array}{l}
(4)\quad f^{(\nu)}(x),\ f^{(\nu+1)}(x),\ f^{(\nu+2)}(x),\ .\ .\ .\ f^{(\nu+\mu-1)}(x),\ f^{\nu+\mu}(x) \\[4pt]
\delta_1,\quad (-1)^\mu \quad -(-1)^\mu \quad (-1)^\mu \quad .\ .\ . \quad - \quad + \\[2pt]
\delta_2,\quad + \quad + \quad + \quad .\ .\ . \quad + \quad +
\end{array}
\] % [sic] last term of row (4) prints as $f^{\nu+\mu}(x)$, without parentheses around the exponent, unlike the other terms in the same row

that is, in the sequence (3), on passing through $\xi$, nothing but changes of sign become successions of signs, and $\mu$ changes of sign are lost in the sequence (3).

If $f^{(\nu+\mu)}(\xi)$ is negative, then all the signs in (4) are the opposite ones, and the conclusion otherwise remains the same.

Now if $\nu = 0$, that is, if $f^{(\nu)}(x)$ is the original function $f(x)$ itself, which then has a $(\mu+1)$-fold root at $\xi$, then on passing through $\xi$ a loss of $\mu$ changes of sign also takes place in the sequence (1).

But if $\nu > 0$ and the function $f^{(\nu-1)}(x)$ preceding $f^{(\nu)}(x)$ is different from zero, then we have the following signs:

1. $\mu$ even:

\[
\begin{array}{ll}
\text{a)}\ f^{(\nu-1)}(x),\ f^{(\nu)}(x), & \text{b)}\ f^{(\nu-1)}(x),\ f^{(\nu)}(x) \\[2pt]
\delta_1,\quad + \quad + & \quad - \quad + \\[2pt]
\delta_2.\quad + \quad + & \quad - \quad +
\end{array}
\]
